\section{A uniform low eigenvalue and a balanced region}
\label{sec:mixed-inertia}

The Schur complement of Proposition~\ref{prop:low-spectrum} also
counts roots when its diagonal has mixed signs. A comparison at three
gives a low eigenvalue for every positive real parameter triple.

\begin{lemma}
\label{lem:mixed-inertia}
For $0<x<a$, suppose all three $k_t(x)$ are nonzero. Let $q$ be the
number of negative diagonal values and put
\[
F(x)=1-\frac{a}{k_a(x)}-\frac{b}{k_b(x)}-\frac{c}{k_c(x)}.
\]
Then $n_-(K(x))=q+\mathbf1_{\{F(x)<0\}}$ and
$\operatorname{nullity}K(x)=\mathbf1_{\{F(x)=0\}}$.
The number of positive complement quotient eigenvalues strictly below
$x$ is $n_-(K(x))-1$.
\end{lemma}

\begin{proof}
The bordered matrix
\[
\begin{pmatrix}\operatorname{diag}(k_t(x))&v\\v^T&1\end{pmatrix}
\]
is congruent both to $K(x)\oplus[1]$ and to
$\operatorname{diag}(k_t(x))\oplus[F(x)]$. Sylvester's law of inertia
gives the formulas. The block $D-xI$ is positive definite, and $C$ has
a simple zero eigenvalue; subtracting this eigenvalue gives the stated count.
\end{proof}

A direct comparison at three avoids the reciprocal formula at zero
diagonal entries and supplies a uniform low positive root.

\begin{theorem}
\label{thm:uniform-low-root}
For every ordered real triple $0<a\leq b\leq c$, the six-support
complement quotient has a nonzero eigenvalue in $(0,3)$.
\end{theorem}

\begin{proof}
Use $s=a+b+c$, $P=abc$ and $v=(\sqrt a,\sqrt b,\sqrt c)^T$ from
Proposition~\ref{prop:low-spectrum}, and define
\[
L=\operatorname{diag}\left(s-\frac{3P}{a^2},
s-\frac{3P}{b^2},s-\frac{3P}{c^2}\right)-vv^T.
\]
Expansion of this diagonal-minus-rank-one determinant gives
\begin{align*}
\det L&=6\bigl(c(a-b)^2+b(a-c)^2+a(b-c)^2\bigr),\\
v^TLv&=-3P\left(\frac1a+\frac1b+\frac1c\right)<0.
\end{align*}
The determinant and Rayleigh value give exactly two negative eigenvalues
at unequal parameters. At equal parameters $L=-vv^T$; thus $L$ is
nonpositive on a two-dimensional subspace in both cases.

For $a>3$, the identity
\[
K(3)=L-\operatorname{diag}_{t\in\{a,b,c\}}
\left(3+\frac{9P}{t^2(t-3)}\right)
\]
subtracts a positive definite matrix. Since $D-3I>0$, this gives
at least two negative eigenvalues of $\mathcal C-3I$.

If $b\leq3$, restrict $\mathcal C-3I$ to the first two singleton
supports and their paired supports. Its blocks are
$\bigl[\begin{smallmatrix}A_0&-\sqrt P I_2\\-\sqrt P I_2&D_0\end{smallmatrix}\bigr]$
with $D_0\leq0$. On vectors $(z,tz)$, the quadratic form is at most
$z^T(A_0-2t\sqrt P I_2)z$, which is negative definite for large $t$.

It remains to consider $a\leq3<b$. Eliminate the positive paired
diagonals $b-3,c-3$. On singleton supports $2,3$, the vector
$(\sqrt b,\sqrt c)$ has quadratic value
\[
 q=(a-3)(b+c)-3P\left(\frac1{b-3}+\frac1{c-3}\right)<0.
\]
Compress the remaining matrix to singleton $1$, its paired support,
and this vector. The resulting matrix is
\[
 \begin{pmatrix}M&-\sqrt P&r\\-\sqrt P&a-3&0\\r&0&q\end{pmatrix},
 \quad M=bc+b+c-3>0,\quad r=-\sqrt a(b+c).
\]
Its determinant $M(a-3)q-Pq-(a-3)r^2$ is positive.
Its positive and negative diagonal entries therefore force exactly
two negative eigenvalues. This includes the singular boundary $a=3$.

In every case $\mathcal C-3I$ has at least two negative eigenvalues.
Since $\mathcal C$ is positive semidefinite with a simple zero,
at least one positive eigenvalue lies strictly below three.
\end{proof}

\begin{corollary}
\label{cor:balanced-interval}
For $4\leq a\leq b\leq c$, if
\begin{equation}
\label{eq:balanced-interval}
(a-2)(a+b+c-2)>2bc,
\end{equation}
then $C$ has an eigenvalue in $(2,3)$, and $G_{p^a q^b r^c}$ is not
Laplacian integral. Nonzero values of both $h_C(1)$ and $h_C(2)$ also
imply nonintegrality. Hence $h_C(1)h_C(2)=0$ is necessary for an integral graph.
\end{corollary}

\begin{proof}
Condition~\eqref{eq:balanced-interval} is exactly $k_a(2)>0$, so all
three diagonal values are positive. As in
Proposition~\ref{prop:low-spectrum}, congruence turns $K(2)$ into
$I-zz^T$, with $\|z\|^2=\sum_t t/k_t(2)>1$, since every $k_t(2)<s$.
It has one negative and two positive eigenvalues and no zero eigenvalue.
Consequently the only quotient eigenvalue at or below two is zero.
The positive root below three from Theorem~\ref{thm:uniform-low-root}
lies in $(2,3)$ and lifts to $(V-3,V-2)$.
For the endpoint assertion, a positive root in $(0,3)$ can be an integer
only if it equals one or two, giving the necessary condition.
\end{proof}

\begin{corollary}
\label{cor:balanced-span}
Let $a$ be the minimum exponent and $\rho$ the maximum minus minimum.
Every positive three-prime exponent triple with
\[
a>\rho+4+\sqrt3(\rho+2)
\]
gives a nonintegral ideal intersection graph. In particular,
$a\geq3\rho+8$ still suffices.
\end{corollary}

\begin{proof}
Write the ordered triple as $(a,a+d,a+\rho)$ with $0\leq d\leq\rho$.
The left side minus the right side of~\eqref{eq:balanced-interval} is
\[
a^2-2a\rho-8a-2\rho^2-4\rho+4+(\rho-d)(a+2\rho+2).
\]
The last term is nonnegative, while the first part factors as
\[
\bigl(a-\rho-4-\sqrt3(\rho+2)\bigr)
\bigl(a-\rho-4+\sqrt3(\rho+2)\bigr).
\]
Both factors are positive under the stated bound, which also gives
$a>4+2\sqrt3>4$. Corollary~\ref{cor:balanced-interval} applies.
The old linear bound is larger, since
$3\rho+8-[\rho+4+\sqrt3(\rho+2)]=(2-\sqrt3)(\rho+2)>0$.
\end{proof}

The new threshold is exact for guaranteeing the strict diagonal
criterion~\eqref{eq:balanced-interval} from the minimum and span alone:
at fixed $a,\rho$ its left side minus right side is minimized by
$b=c=a+\rho$, and it vanishes at the displayed threshold.
This optimality concerns that sufficient criterion, not the entire
spectral region where a root lies in $(2,3)$. For example,
$(35,45,45)$ and $(281,381,381)$ satisfy the new bound but not
$a\geq3\rho+8$.

For example, $(20,21,24)$, $(38,43,48)$ and $(308,350,408)$ satisfy
the span criterion and lie outside the four fixed-gap families of
Theorem~\ref{thm:bounded-gaps}. The argument allows an unbounded span
and uses no exponent scan or modular endpoint certificate. The necessary
endpoint-zero condition is not sufficient for integrality: those
triples may still have other noninteger quotient eigenvalues.
